\chapter{Regular Expressions}
\label{sec:regular-expressions}

%

\section{Syntax}
\label{sec:regex-syntax}

Let $\Sigma$ be an alphabet. Regular expressions over $\Sigma$ are
defined inductively by the grammar:
\begin{equation}\label{eq:regex-grammar}
r,s ::= x \mid \varepsilon \mid \emptyset
      \mid r+s \mid r\cdot s \mid r^*
\qquad (x\in\Sigma).
\end{equation}

In the grammar, $r$ and $s$ range over regular expressions, while $x$
ranges over the alphabet $\Sigma$ and denotes an arbitrary symbol, such as
$a$, $b$, or $c$. The constants $\varepsilon$ and $\emptyset$ represent the 
\emph{empty-word} and \emph{empty-language} expressions. The operators $+$, $\cdot$,
and $^*$ denote alternation, concatenation, and Kleene star,
respectively.

Concatenation is written implicitly, so $rs$ denotes $r \cdot s$.
Given this notation, we adopt the precedence and associativity
conventions. Unless parentheses are shown explicitly, Kleene star binds
most tightly, followed by concatenation, followed by alternation. For example,
\begin{align*}
ab^*c &= (a\cdot(b^*))\cdot c, \\
a+bc^* &= a+(b\cdot(c^*)).
\end{align*}
For repeated operators we use left association, so
\begin{align*}
a+b+c &= (a+b)+c, \\
abc &= (ab)c.
\end{align*}

The grammar above is the core syntax. Most regular expression libraries 
also offer one-or-more, optional, and bounded repetition ($r^+$, $r?$, 
$r^{\{m,n\}}$), but these add no expressive power since each can be defined 
by the core constructs: $r^+ = r\cdot r^*$, $r? = r+\varepsilon$, and
\[
r^{\{m,n\}} = \underbrace{r\cdot\ldots\cdot r}_{m\text{ times}}
\cdot
\underbrace{(r+\varepsilon)\cdot\ldots\cdot(r+\varepsilon)}_{n-m\text{ times}}.
\]

In the implementation, expressions using such operators are translated
into the core constructs before they reach the parsing algorithms.

As a running example, consider the regular expression over
$\Sigma=\{0,1\}$
\begin{equation}\label{eq:r01-def}
r_{01} = (0+1)^*\cdot0\cdot1.
\end{equation}

By the conventions above, its syntax tree is
$((0+1)^*\cdot0)\cdot1$. This tree represents the hierarchical structure
of the expression: a Kleene star applied to the alternation $0+1$, followed by
concatenation with the literals $0$ and $1$.

The syntactic structure must be given a semantic interpretation.

%

\section{Semantics}
\label{sec:regex-semantics}

A semantic interpretation assigns a language to each regular expression. A
\emph{word} over $\Sigma$ is a finite sequence of symbols from
$\Sigma$. The set of all words over $\Sigma$, including the empty word
$\varepsilon$, is denoted by $\Sigma^*$. A \emph{language} over
$\Sigma$ is a subset of $\Sigma^*$.

For a regular expression $r$, we write $L(r)$ for the language it denotes,
where $L(r) \subseteq \Sigma^*$. The definition follows the structure of the syntax tree by
\emph{structural recursion}. Each case corresponds to one production of
the grammar in \eqref{eq:regex-grammar}:
\begin{equation}\label{eq:semantics}
\begin{aligned}
L(x) &= \{x\}, \\
L(\varepsilon) &= \{\varepsilon\}, \\
L(\emptyset) &= \emptyset, \\
L(r+s) &= L(r)\cup L(s), \\
L(r\cdot s) &= \{vw \mid v\in L(r),\ w\in L(s)\}, \\
L(r^*) &= \{\varepsilon\}
          \cup
          \{w_1\cdots w_n
            \mid n\geq 1,\ w_1,\ldots,w_n\in L(r)\}.
\end{aligned}
\end{equation}

Table~\ref{tab:lang-examples} applies these definitions to small expressions. 
The literal $a$ denotes the language containing only the word $a$. The expression 
$\varepsilon$ denotes the language containing only the empty word, while $\emptyset$ 
denotes the language containing no words. The alternation $a+b$ combines the languages of $a$ and 
$b$ to give $\{a,b\}$, and concatenation $ab$ combines words from the languages of 
$a$ and $b$ to give $\{ab\}$. 
The Kleene star $a^*$ allows zero or more concatenations of words from
the language of $a$, giving $\{\varepsilon,a,aa,aaa,\ldots\}$.
\begin{table}[ht]
\centering
\begin{tabular}{cc}
\hline
Expression $r$ & Language $L(r)$\\
\hline
$a$ & $\{a\}$ \\
$\varepsilon$ & $\{\varepsilon\}$ \\
$\emptyset$ & $\emptyset$ \\
$a+b$ & $\{a,b\}$ \\
$ab$ & $\{ab\}$ \\
$a^*$ & $\{\varepsilon,a,aa,aaa,\ldots\}$ \\
\hline
\end{tabular}
\caption{Languages denoted by basic expressions.}
\label{tab:lang-examples}
\end{table}

For example $r_{01} = (0+1)^*\cdot0\cdot1$ from \eqref{eq:r01-def}, we compute 
its language step by step:
\begin{equation}\label{eq:r01-language}
\begin{aligned}
L(0) &= \{0\}, \\
L(1) &= \{1\}, \\
L(0+1) &= L(0) \cup L(1) = \{0,1\}, \\
L((0+1)^*) &= \{0,1\}^*, \\
L(r_{01}) &= L((0+1)^*) \cdot L(0) \cdot L(1) = \{\, w \cdot 0 \cdot 1 \mid w\in\{0,1\}^* \,\}.
\end{aligned}
\end{equation}

This is the set of binary words ending in $01$. For example, the words
$01$, $001$, and $1101$ belong to $L(r_{01})$, whereas $\varepsilon$,
$0$, and $010$ do not.

This is the question of whether a word \emph{matches} a regular
expression: $w$ \emph{matches} $r$ if $w\in L(r)$. Determining this is
the \emph{membership problem}.

%

\section{Rust Representation}
\label{sec:regex-rust-implementation}
The grammar from ~\eqref{eq:regex-grammar} is represented 
directly as a Rust enum with one variant for 
each syntactic construct, as shown in Figure~\ref{fig:regex-enum}:
\begin{figure}[ht]
\small
\begin{verbatim}
#[derive(Debug, Clone, PartialEq, Eq, Hash)]
pub enum Regex {
    Phi, Eps, Lit(char),
    Seq(Box<Regex>, Box<Regex>),  // r . s
    Alt(Box<Regex>, Box<Regex>),  // r + s
    Star(Box<Regex>),              // r*
}

impl Regex {
    pub fn seq(r: Regex, s: Regex) -> Regex {
        Regex::Seq(Box::new(r), Box::new(s))
    }
    pub fn alt(r: Regex, s: Regex) -> Regex {
        Regex::Alt(Box::new(r), Box::new(s))
    }
    pub fn star(r: Regex) -> Regex { Regex::Star(Box::new(r)) }
    pub fn lit(c: char) -> Regex { Regex::Lit(c) }
}
\end{verbatim}
\caption{The \texttt{Regex} enum and constructors. Derived traits support
operations needed by parsing: \texttt{Clone} for duplication,
\texttt{PartialEq}/\texttt{Eq} for comparison, \texttt{Hash} for collections.}
\label{fig:regex-enum}
\end{figure}

The type carries no precedence or associativity information, as these are
determined by the syntax tree. Rust requires explicit indirection because
enum sizes must be known at compile time. Recursive variants therefore use
\texttt{Box<Regex>}. A direct definition such as
\texttt{Seq(Regex, Regex)} would have infinite size, whereas \texttt{Box}
stores recursive values separately, providing finite size
~\citep{rust-book-box,rust-reference-recursive-types}. This differs from
Haskell's heap-allocated data types.

The running example $r_{01}=(0+1)^*\cdot0\cdot1$ from ~\eqref{eq:r01-def}
is represented by the syntax tree:
\begin{equation}\label{eq:r01-tree}
\operatorname{Seq}\bigl(
    \operatorname{Seq}(
        \operatorname{Star}(\operatorname{Alt}(0,1)),
        0
    ),
    1
\bigr).
\end{equation}
In Rust, this is represented as shown in Figure~\ref{fig:r01-rust}, with the
outermost \texttt{Seq} containing \texttt{Lit('1')} and 
\texttt{Seq(Star(Alt(0,1)), 0)} as its right and left children.
\begin{figure}[ht]
\small
\begin{verbatim}
Regex::seq(
    Regex::seq(
        Regex::star(Regex::alt(Regex::lit('0'), Regex::lit('1'))),
        Regex::lit('0')),
    Regex::lit('1'))
\end{verbatim}
\caption{Rust representation of $r_{01}$.}
\label{fig:r01-rust}
\end{figure}